\section{Dinámica de filtrado probabilístico entre capas}

El análisis empírico revela una dinámica de \textbf{filtrado probabilístico bidireccional} entre niveles de integración de la persistencia. Este filtrado es un \textbf{efecto estructural} que emerge de la configuración de acoplamiento entre componentes: no se postula ninguna fuerza activa, agente o intención. Se describe, de manera estrictamente operacional, una variación observable y medible en la probabilidad de que una intensificación local (Capa 1) genere reorganización estructural (Capa 3) según el nivel de antisincronización del sistema. A este efecto lo denominamos, entre comillas y con valor heurístico secundario, «resistencia ascendente» o «resistencia descendente», según la dirección de la jerarquía. El concepto primario y riguroso es el de filtrado probabilístico; los términos entre comillas sirven como etiquetas intuitivas pero subordinadas.

\subsection{Filtrado probabilístico ascendente (o «resistencia ascendente»)}

Cuando la antisincronización es elevada, la probabilidad de que una intensificación local (Capa 1) logre modificar la matriz de coherencias del sistema (Capa 3) disminuye sustancialmente. Los módulos mantienen escalas de persistencia $\tau_s$ dispersas; la fragmentación resultante reduce la permeabilidad relacional necesaria para que los cambios locales se propaguen. En términos estrictamente operacionales, el aumento de la desviación estándar de $\tau_s$ entre componentes actúa como un filtro que eleva el umbral requerido para que una intensificación local se traduzca en reorganización estructural detectable mediante la norma de Frobenius. La evidencia de la Sección 3 ilustra directamente este mecanismo: los picos hiperestructurales ocurren preferentemente tras reducción de antisincronización (media significativamente inferior en ventanas de $\pm 50$ pasos, $p=0.0229$ en el mapa logístico; valores aún más bajos en los 100 pasos previos). La reducción de fragmentación eleva la probabilidad de transición de Capa 1 a Capa 3. Este es el filtrado probabilístico en dirección ascendente.

\subsection{Filtrado probabilístico descendente (o «resistencia descendente»)}

Cuando la antisincronización ha descendido (Capa 2 consolidada), la probabilidad de que intensificaciones locales aún fragmentadas alteren la estructura global ya alcanzada también se reduce. Una organización coherente de escalas de persistencia genera una configuración en la que las fluctuaciones locales que no alcanzan un umbral mínimo de reducción de fragmentación tienden a ser absorbidas sin modificar la matriz de relaciones global. Operacionalmente, la baja antisincronización estabiliza la estructura global frente a intensificaciones locales que no logran alinear suficientemente las escalas de persistencia de los módulos. La robustez al ruido documentada en la Sección 3 (descenso de 121 a 100 picos con 0--15 \% de ruido gaussiano, pero invariancia esencial de los cambios por norma de Frobenius) apoya este filtrado descendente: una vez alcanzada la coherencia de escalas, la estructura global (Capa 3) tiende a absorber fluctuaciones locales fragmentadas sin reorganizarse. Este es el filtrado probabilístico en dirección descendente.

\subsection{Implicación ontológica}

La antisincronización funciona como regulador de la permeabilidad entre capas: su valor modula la probabilidad de transición de la Capa 1 a la Capa 3 mediante el efecto estructural de filtrado probabilístico bidireccional. El presente de nivel superior (Capa 3) no se reduce a la suma de presentes locales intensificados; emerge como una estructura relacional cuya estabilidad depende de la reducción previa de fragmentación. Las transiciones entre capas siguen siendo los eventos discretos ontológicamente relevantes del tiempo extramental detectable mediante RECD. (Los términos «resistencia ascendente» y «resistencia descendente» se usan solo ocasionalmente como etiquetas heurísticas secundarias para los dos sentidos del filtrado probabilístico; no designan procesos activos ni protectores.)